\clearpage
\section{Chứng minh đọc sâu}
\label{sec:proofs}
Phần này có thể đọc ở lượt thứ hai. Các chứng minh làm rõ vì sao công
thức có dạng như trong mạch chính; không cần thuộc lòng từng dòng để bắt
đầu giải bài. Những định lý nền như CLT và phân phối Hotelling được dùng
như kết quả đã có trong giáo trình, không tuyên bố chứng minh đầy đủ ở đây.

\subsection{Làm trắng và phân phối khoảng cách}
\label{app:whitening}
Vì $\Sigma$ đối xứng xác định dương, tồn tại phân rã phổ
\[
 \Sigma=QDQ',\qquad Q'Q=I,
 \quad D=\diag(\lambda_1,\ldots,\lambda_p),\quad\lambda_i>0.
\]
Đặt $\Sigma^{-1/2}=QD^{-1/2}Q'$ và
$\mathbf Z=\Sigma^{-1/2}(\mathbf X-\vmu)$.
Theo tính chất biến đổi tuyến tính,
\[
 \mathbf Z\sim\N_p(0,\Sigma^{-1/2}\Sigma\Sigma^{-1/2})=\N_p(0,I).
\]
Các thành phần của $\mathbf Z$ chuẩn đồng thời và không tương quan, nên
độc lập. Mặt khác,
\[
 (\mathbf X-\vmu)'\Sigma^{-1}(\mathbf X-\vmu)
 =\mathbf Z'\mathbf Z=\sum_{i=1}^p Z_i^2\sim\chi^2_p.
\]
Đây là chứng minh Result 4.7(a). Việc lấy phân vị của phân phối vừa tìm
cho ngay kết luận về xác suất của elipsoid.
\nguon{Chương 4, §4.2, Result 4.7 và chứng minh, tr.~163--164}

\subsection{Phân phối có điều kiện qua phần dư}
Đặt $B=\Sigma_{12}\Sigma_{22}^{-1}$ và
\[
 \mathbf R=\mathbf X_1-\vmu_1-B(\mathbf X_2-\vmu_2).
\]
Vectơ $(\mathbf R',\mathbf X_2')'$ chuẩn vì là biến đổi tuyến tính của
vectơ gốc. Ta tính
\[
 \Cov(\mathbf R,\mathbf X_2)=\Sigma_{12}-B\Sigma_{22}=0.
\]
Do đó $\mathbf R$ độc lập với $\mathbf X_2$.
Hiệp phương sai phần dư là
\begin{align*}
 \Cov(\mathbf R)
 &=\Sigma_{11}-B\Sigma_{21}-\Sigma_{12}B'+B\Sigma_{22}B'\\
 &=\Sigma_{11}-\Sigma_{12}\Sigma_{22}^{-1}\Sigma_{21}.
\end{align*}
Ở dòng cuối ta dùng tính đối xứng của $\Sigma_{22}$.
Khi biết $\mathbf X_2=\mathbf x_2$,
\[
 \mathbf X_1=\vmu_1+B(\mathbf x_2-\vmu_2)+\mathbf R.
\]
Phần trước $\mathbf R$ là hằng số có điều kiện, còn $\mathbf R$ vẫn giữ
nguyên phân phối. Suy ra cả trung bình và hiệp phương sai trong
\eqref{eq:conditionalmean}--\eqref{eq:conditionalcov}.
\nguon{Chương 4, §4.2, chứng minh Result 4.6, tr.~161}

Hiệp phương sai được giảm theo nghĩa ma trận:
\[
 \Sigma_{11}-\Sigma_{11\mid2}
 =\Sigma_{12}\Sigma_{22}^{-1}\Sigma_{21}\succeq0.
\]
Thật vậy với mọi $\mathbf a$,
$\mathbf a'\Sigma_{12}\Sigma_{22}^{-1}\Sigma_{21}\mathbf a\ge0$.
Do đó phương sai của \emph{bất kỳ tổ hợp tuyến tính} của khối mục tiêu
không tăng khi điều kiện hóa trong mô hình chuẩn này. ``Giảm theo nghĩa
ma trận'' không có nghĩa mọi phần tử hiệp phương sai đều giảm về mặt số học.

\subsection{Vì sao trung bình mẫu độc lập với \texorpdfstring{$S$}{S}?}
\label{app:wishart}
Viết bảng quan sát thành ma trận $D$ kích thước $n\times p$.
Chọn ma trận trực giao $H$ kích thước $n\times n$, có hàng đầu
$\mathbf 1'/\sqrt n$, các hàng còn lại vuông góc với $\mathbf1$.
Đặt $Y=HD$, với các hàng được xem là vectơ quan sát mới.

Hàng đầu là $\sqrt n\Xbar'$ và có trung bình $\sqrt n\vmu'$.
Các hàng còn lại có trung bình 0. Hiệp phương sai giữa hàng $r$ và hàng
$s$ của $Y$ là
\[
 \sum_{j=1}^n H_{rj}H_{sj}\Sigma
 =\begin{cases}\Sigma,&r=s,\\0,&r\ne s.\end{cases}
\]
Các hàng chuẩn đồng thời, nên độc lập. Viết $\mathbf Y_r$ là chuyển vị
của hàng $r$, ta có $\mathbf Y_2,\ldots,\mathbf Y_n$ iid $\N_p(0,\Sigma)$.
Tính trực giao cho
\[
 (n-1)S=D'\left(I_n-\frac{\mathbf1\mathbf1'}n\right)D
 =\sum_{r=2}^n\mathbf Y_r\mathbf Y_r'.
\]
Vế phải là Wishart với $n-1$ bậc tự do, chỉ phụ thuộc các hàng 2 đến $n$.
Hàng đầu quyết định trung bình mẫu và độc lập với những hàng đó, nên
$\Xbar\perp S$.
\nguon{Khai triển bổ sung cho kết quả phân phối mẫu (4-23), Chương 4, §4.4, tr.~174}

\subsection{Kiểm tra cực đại MLE bằng trị riêng}
\label{app:mle_global}
Nếu muốn tránh đạo hàm ma trận ở Phần~\ref{sec:mle_derivation}, có thể
kiểm tra cực đại toàn cục bằng bất đẳng thức một biến.
Đặt $M=A/n\succ0$ và $B=M^{1/2}\Sigma^{-1}M^{1/2}\succ0$.
Gọi $\lambda_1,\ldots,\lambda_p$ là các trị riêng của $B$.
Từ log-hợp lý đã đặt $\vmu=\xbar$, ta có
\begin{align*}
 \ell(\xbar,\Sigma)-\ell(\xbar,M)
 &=\frac n2\bigl[\log|B|-\tr(B)+p\bigr]\\
 &=\frac n2\sum_{i=1}^p(\log\lambda_i-\lambda_i+1)\le0.
\end{align*}
Dòng cuối dùng $\log u\le u-1$ với $u>0$, dấu bằng chỉ khi $u=1$.
Cực đại đạt duy nhất khi mọi $\lambda_i=1$, tức $B=I$ vì $B$ đối xứng,
và do đó $\Sigma=M=A/n$. Bước tìm trung bình đã bảo đảm
$\ell(\vmu,\Sigma)\le\ell(\xbar,\Sigma)$, nên đây là cực đại đồng thời.
\nguon{Lập luận tương đương Chương 4, §4.3, Result 4.10, tr.~170--171,
 và Result 4.11, tr.~171--172}

Nếu $A$ có hạng $r<p$, viết
$A=U\diag(a_1,\ldots,a_r,0,\ldots,0)U'$ với $a_i>0$.
Chọn
\[
 \Sigma_\varepsilon=U\diag(a_1/n,\ldots,a_r/n,
 \varepsilon,\ldots,\varepsilon)U',\qquad\varepsilon>0.
\]
Khi đó $\tr(\Sigma_\varepsilon^{-1}A)=nr$ không đổi, nhưng
$\log|\Sigma_\varepsilon|=\text{hằng số}+(p-r)\log\varepsilon
\to-\infty$ khi $\varepsilon\downarrow0$.
Suy ra $\ell(\xbar,\Sigma_\varepsilon)\to+\infty$.
Không có cực đại hữu hạn trong miền hiệp phương sai xác định dương;
ma trận suy biến $A/n$ không khắc phục vấn đề bằng cách được gọi là
nghiệm trong mô hình không suy biến.

Chứng minh bất biến của $T^2$, gồm cả giá trị p và quyết định kiểm định,
được trình bày ngay trong Phần~\ref{sec:invariance_proof} để thuận tiện
đọc cùng ví dụ. Không cần lặp lại chứng minh đó ở đây.

\subsection{Tỷ số định thức và Wilks' lambda}
Đặt $A=(n-1)S$ và $\mathbf r=\xbar-\vmu_0$. Tách tổng bình phương quanh
giá trị giả thuyết:
\begin{align*}
 \sum_j(\mathbf x_j-\vmu_0)(\mathbf x_j-\vmu_0)'
 &=\sum_j\bigl[(\mathbf x_j-\xbar)+\mathbf r\bigr]
 \bigl[(\mathbf x_j-\xbar)+\mathbf r\bigr]'\\
 &=A+n\mathbf r\mathbf r'.
\end{align*}
Hai hạng chéo bằng 0 vì $\sum_j(\mathbf x_j-\xbar)=0$.
Bổ đề định thức cho cập nhật hạng một là
\[
 |A+\mathbf u\mathbf v'|=|A|(1+\mathbf v'A^{-1}\mathbf u).
\]
Lấy $\mathbf u=n\mathbf r$, $\mathbf v=\mathbf r$:
\[
 \frac{|A+n\mathbf r\mathbf r'|}{|A|}
 =1+n\mathbf r'A^{-1}\mathbf r=1+\frac{T^2}{n-1}.
\]
Vì $n\widehat\Sigma=A$ và $n\widehat\Sigma_0=A+n\mathbf r\mathbf r'$,
ta có
\[
 \frac{|\widehat\Sigma|}{|\widehat\Sigma_0|}
 =\left(1+\frac{T^2}{n-1}\right)^{-1}.
\]
Để hoàn tất kết nối với hợp lý, ở MLE của mô hình chuẩn, phần mũ của hàm
hợp lý bằng $e^{-np/2}$ ở cả hai mô hình. Hệ số $(2\pi)^{-np/2}$ cũng
giống nhau, nên tỷ số cực đại chỉ còn
$\Lambda_{\mathrm{LR}}=(|\widehat\Sigma|/|\widehat\Sigma_0|)^{n/2}$.
\nguon{Chương 5, §5.3, chứng minh Result 5.1, tr.~218; phân rã tổng bình phương (4-14), tr.~169}

\subsection{Hình chiếu elipsoid và khoảng đồng thời}
\label{app:projection}
Cho $\mathbf h=\vmu-\xbar$. Nếu $\vmu\in\mathcal C$, thì
$\mathbf h'S^{-1}\mathbf h\le c_T^2/n$.
Viết
\[
 \mathbf a'\mathbf h=(S^{1/2}\mathbf a)'(S^{-1/2}\mathbf h).
\]
Bất đẳng thức Cauchy--Schwarz cho
\begin{align*}
 |\mathbf a'\mathbf h|
 &\le\sqrt{\mathbf a'S\mathbf a}
 \sqrt{\mathbf h'S^{-1}\mathbf h}\\
 &\le\sqrt{c_T^2}\sqrt{\mathbf a'S\mathbf a/n}.
\end{align*}
Đây chính là hai đầu khoảng \eqref{eq:t2interval}. Với $\mathbf a\ne0$,
dấu bằng đạt tại
\[
 \mathbf h=\pm\sqrt{\frac{c_T^2}{n}}
 \frac{S\mathbf a}{\sqrt{\mathbf a'S\mathbf a}}.
\]
Vì thế khoảng là hình chiếu \emph{chính xác} của elipsoid, không phải
một xấp xỉ về hình học.

Theo chiều ngược lại, nếu khoảng phủ cho mọi $\mathbf a$, chọn
$\mathbf a=S^{-1}\mathbf h$ để suy ra
$\mathbf h'S^{-1}\mathbf h\le c_T^2/n$ (trường hợp $\mathbf h=0$ hiển nhiên).
Do đó biến cố ``mọi khoảng đều phủ'' tương đương ``elipsoid phủ vectơ
trung bình'', có xác suất $1-\alpha$ dưới mô hình chuẩn.
\nguon{Chương 5, §5.4, (5-23) và Result 5.3, tr.~225; Supplement 5A về hình chiếu}

\subsection{Bonferroni và Slutsky: hai kiểu bảo đảm khác nhau}
Cho $C_j$ là biến cố khoảng thứ $j$ phủ đúng. Với mọi kết quả ngẫu nhiên,
\[
 \mathbf1_{\cup_j C_j^c}\le\sum_j\mathbf1_{C_j^c}.
\]
Lấy kỳ vọng và phần bù:
\[
 \Prob\left(\bigcap_j C_j\right)
 =1-\Prob\left(\bigcup_j C_j^c\right)
 \ge1-\sum_j\Prob(C_j^c).
\]
Nếu mỗi khoảng có xác suất sai $\alpha/m$, cả nhóm có xác suất phủ
ít nhất $1-\alpha$. Bất đẳng thức luôn đúng; việc xác suất sai của từng
khoảng có \emph{đúng} bằng $\alpha/m$ hay chỉ xấp xỉ tùy vào phương pháp
dựng khoảng và mô hình.
\nguon{Chương 5, §5.4, (5-28), tr.~232}

Với Slutsky, đặt $\mathbf Z_n=\sqrt n(\Xbar-\vmu)$.
CLT cho $\mathbf Z_n\xrightarrow{d}\mathbf Z\sim\N_p(0,\Sigma)$.
Luật số lớn và tính liên tục của phép nghịch đảo tại ma trận xác định
dương cho $S^{-1}\xrightarrow{P}\Sigma^{-1}$.
Slutsky và phép ánh xạ liên tục dẫn đến
\[
 \mathbf Z_n'S^{-1}\mathbf Z_n\xrightarrow{d}
 \mathbf Z'\Sigma^{-1}\mathbf Z\sim\chi^2_p.
\]
Không cần $\Xbar$ và $S$ độc lập ở mẫu hữu hạn cho lập luận giới hạn này.
Đây là một lý do suy diễn mẫu lớn có thể vượt ra ngoài tổng thể chuẩn.
\nguon{Diễn giải bổ sung cho Chương 4, §4.5, (4-28), tr.~176}

\clearpage
\section{Đáp án và lời giải ngắn cho bài tự luyện}
\label{app:answers}
\subsection*{Bài 1--3: đọc tham số và biến đổi}
\noindent\textbf{Bài 1.}
Trung bình là 10 và 20; độ lệch chuẩn là 2 và 3; tương quan
$2/(2\cdot3)=1/3$. $|\Sigma|=36-4=32>0$, $\sigma_{11}>0$, nên
xác định dương. Nghịch đảo là
$\frac1{32}\begin{pmatrix}9&-2\\-2&4\end{pmatrix}$.
Các phương sai không cho biết hiệp phương sai ngoài đường chéo; thiếu
nó, hướng và độ dài các trục elip chưa được xác định.

\noindent\textbf{Bài 2.}
Ma trận chọn $A=\begin{pmatrix}0&0&1\\1&0&0\end{pmatrix}$, nên
\[
 (X_3,X_1)'\sim\N_2\left((2,1)',
 \begin{pmatrix}2&-1\\-1&4\end{pmatrix}\right),\qquad
 X_2\sim\N(-1,5).
\]
Không biết giá trị biến nào khác, nên đây là phân phối biên.

\noindent\textbf{Bài 3.}
$\mathbf a=(1,2,-1)'$, $\mathbf a'\vmu=-3$ và
$\Sigma\mathbf a=(5,10,-3)'$, do đó phương sai $5+20+3=28$.
Cách khai triển: $4+4\cdot5+2+2(1)(-1)(-1)=28$.
Hạng chéo là dương vì hệ số của $X_3$ âm và hiệp phương sai với $X_1$
cũng âm. Vậy $Y\sim\N(-3,28)$.

\subsection*{Bài 4--6: độc lập và điều kiện}
\noindent\textbf{Bài 4.}
$\Cov(U,V)=\Var(X_1)-\Var(X_3)=2\ne0$: không độc lập.
$\Cov(X_2,U)=0+0=0$: độc lập nhờ chuẩn đồng thời.
Với $U_c=X_1+cX_3$,
\[
 \Cov(U_c,V)=4-(-1)+c(-1-2)=5-3c.
\]
Nghiệm $c=5/3$.

\noindent\textbf{Bài 5.}
$\Cov(Z,Z^2)=\E(Z^3)=0$, nhưng $Z^2$ được xác định bởi $Z$, nên phụ
thuộc. Để minh họa chuẩn của từng biên vẫn chưa đủ, dùng $V=BZ$ với $B$
độc lập và nhận $\pm1$ đồng xác suất. Hai biên chuẩn, hiệp phương sai
0, nhưng $|V|=|Z|$; chúng không chuẩn đồng thời và không độc lập.

\noindent\textbf{Bài 6.}
Khối mục tiêu $\mathbf U=(X_1,X_2)'$ có
$\vmu_U=(1,-1)'$, $\Sigma_{UU}=\diag(4,5)$.
Khối đã biết là $V=X_3$, $\mu_V=2$, $\Sigma_{VV}=2$;
$\Sigma_{UV}=(-1,0)'$, $\Sigma_{VU}=(-1,0)$.
\[
 \mathbf U\mid V=4\sim\N_2\left((0,-1)',\diag(7/2,5)\right).
\]
Chỉ thành phần thứ nhất được cập nhật; $X_2$ độc lập với $X_3$.
Trong phân phối có điều kiện này, hai thành phần vẫn độc lập vì ma trận
hiệp phương sai chéo bằng không và chúng chuẩn đồng thời.

\subsection*{Bài 7--9: phân phối mẫu và kiểm định}
\noindent\textbf{Bài 7.}
\[
 \Xbar\sim\N_2\left((10,20)',
 \begin{pmatrix}.16&.12\\.12&.36\end{pmatrix}\right),\qquad
 24S\sim W_{2,24}(\Sigma),\quad\Xbar\perp S.
\]
Độ lệch chuẩn của $X_{j1}$ là 2, của $\bar X_1$ là 0.4.
Với mẫu iid không chuẩn, $\E(\Xbar)=\vmu$, $\Cov(\Xbar)=\Sigma/25$
và $\E(S)=\Sigma$ vẫn đúng khi có moment hữu hạn. Phân phối chuẩn chính
xác, Wishart và độc lập của $\Xbar,S$ không còn được bảo đảm.

\noindent\textbf{Bài 8.}
$n=5,p=2$, $\xbar=(3,3)'$.
\[
 S=\begin{pmatrix}2.5&2\\2&2.5\end{pmatrix},\quad
 |S|=2.25,\quad
 S^{-1}=\begin{pmatrix}10/9&-8/9\\-8/9&10/9\end{pmatrix}.
\]
Với $\mathbf d=(1,1)'$, $T^2=5\cdot4/9=20/9$.
Thống kê $F=(3/8)T^2=5/6\sim F_{2,3}$ dưới $H_0$.
Ngưỡng cho $T^2$ là $(8/3)9.5521\approx25.4723$.
Chưa bác bỏ giả thuyết ở mức 5\%; giá trị p khoảng 0.5154.
Không coi đây là bằng chứng hai trung bình chắc chắn bằng 2.
MLE trung bình là $(3,3)'$; MLE hiệp phương sai là
$\widehat\Sigma=\frac45S=\begin{pmatrix}2&1.6\\1.6&2\end{pmatrix}$.
Trong log-hợp lý, phân rã tổng tích ngoài thành $4S$ và
$5(\xbar-\vmu)(\xbar-\vmu)'$ cho thấy phần phụ thuộc trung bình
đạt cực đại tại $\vmu=\xbar$. Tối ưu phần còn lại cho $\Sigma=4S/5$;
đây là thiết lập ở Phần~\ref{sec:mle_derivation}, không phải quy tắc
đổi mẫu số tùy ý.

\noindent\textbf{Bài 9.}
\[
 \bar{\mathbf y}=(305,-1.7)',\qquad\vmu_{Y,0}=(205,-1.8)',
 \qquad S_Y=\begin{pmatrix}25000&20\\20&.025\end{pmatrix}.
\]
Độ lệch mới là $(100,0.1)'$, $T_Y^2=20/9=T_X^2$.
Trong công thức tổng quát cho bảng hàng, $D_Y=D_XC'+\mathbf1\mathbf d'$;
với $C$ chéo riêng bài này, $C'=C$, nhưng không được suy như vậy cho mọi ma trận.
Vì $n=5,p=2$ và $T^2$ giữ nguyên, thống kê $F$, số bậc tự do và giá trị
p khoảng 0.5154 giữ nguyên; vẫn chưa bác bỏ giả thuyết mức 5\%.
Đẳng thức đại số không cần chuẩn; chuẩn cần cho phân phối $F$ chính xác.

\subsection*{Bài 10--12: khoảng, hợp lý và mẫu lớn}
\noindent\textbf{Bài 10.}
$m=4$. Tổ hợp $\mathbf a=(2,-1)'$ có ước lượng $2(10)-20=0$,
\[
 \mathbf a'S\mathbf a=4(4)+9-4(3)=13,\qquad
 \mathrm{SE}=\sqrt{13/40}\approx0.5701.
\]
Khoảng $T^2$ là $0\pm\sqrt{6.6604}\sqrt{13/40}\approx0\pm1.4713$.
Khoảng Bonferroni là $0\pm t_{39;0.99375}\sqrt{13/40}
\approx0\pm1.4930$.
Với danh sách này, Bonferroni hơi rộng hơn $T^2$; không mặc định nó luôn
hẹp hơn. Không dùng $m=3$ để bảo vệ cả bốn phát biểu.
Khi \emph{chỉ xét hai trung bình}, dùng $m=2$, hình chữ nhật là
$[9.2628;10.7372]\times[18.8942;21.1058]$.
$P$ ở trong chữ nhật nhưng $T^2(P)=19.6>6.6604$, ngoài elip;
$R$ ở ngoài chữ nhật nhưng $T^2(R)=6.4<6.6604$, trong elip.
Đổi dấu hiệp phương sai chéo giữ nguyên chữ nhật thành phần nhưng đổi
hướng elip. Phương sai tổ hợp $(2,-1)'$ tăng từ 13 thành
$16+9-4(-3)=37$, nên sai số chuẩn tăng từ $\sqrt{13/40}$
thành $\sqrt{37/40}$.

\noindent\textbf{Bài 11.}
\[
 \widehat\Sigma=\frac45 S=
 \begin{pmatrix}2&1.6\\1.6&2\end{pmatrix},\quad
 \widehat\Sigma_0=\widehat\Sigma+\mathbf d\mathbf d'
 =\begin{pmatrix}3&2.6\\2.6&3\end{pmatrix}.
\]
Hai định thức là 1.44 và 2.24.
\[
 \lambda_W=\frac9{14},\qquad
 \Lambda_{\mathrm{LR}}=\left(\frac9{14}\right)^{5/2},\qquad
 T^2=4\left(\frac{14}9-1\right)=\frac{20}9.
\]
Giá trị p được tính từ \emph{phân phối dưới $H_0$} của thống kê;
tỷ số định thức không phải phép tính đuôi xác suất nói chung.

\noindent\textbf{Bài 12.}
$T^2=100(0.2^2+0.1^2)=5<5.9915$: chưa bác bỏ ở mức xấp xỉ 5\%.
Giá trị p xấp xỉ $\exp(-5/2)\approx0.0821$.
Nửa độ rộng khoảng chi-bình phương là
$\sqrt{5.9915/100}\approx0.2448$, nên các khoảng là
$[-0.0448;0.4448]$ và $[-0.1448;0.3448]$.
Bonferroni cho $m=2$ có nửa độ rộng 0.2241, cho các khoảng
$[-0.0241;0.4241]$ và $[-0.1241;0.3241]$.
Đây đều là bảo đảm xấp xỉ, không chính xác như mẫu chuẩn hữu hạn.

Với $n=20,p=25$, hạng $S\le19<25$, nên không thể dùng nghịch đảo.
Với $n=2000,p=1900$, việc $n-p=100$ không đủ để khẳng định xấp xỉ cổ
điển tốt: $p$ rất lớn so với $n$, và giới hạn $p$ cố định không tự bảo
đảm cho tình huống cao chiều đó.

\clearpage
\section{Bảng công thức và điều kiện đi kèm}
\label{sec:formula}
\subsection{Công thức phân phối}
\begin{longtable}{L{.28\linewidth}L{.63\linewidth}}
\toprule
Điều kiện / câu hỏi & Công thức\\
\midrule
\endhead
Chuẩn không suy biến &
$f(\mathbf x)=(2\pi)^{-p/2}|\Sigma|^{-1/2}
 \exp[-(\mathbf x-\vmu)'\Sigma^{-1}(\mathbf x-\vmu)/2]$;
$\Sigma$ xác định dương.\\[4pt]
Biến đổi chuẩn & $A\mathbf X+\mathbf d\sim\N_q(A\vmu+\mathbf d,A\Sigma A')$.
Thiếu hạng hàng của $A$ có thể cho chuẩn suy biến.\\[4pt]
Phân phối biên chuẩn & Chọn các thành phần của $\vmu$ và các hàng/cột cùng thứ tự của $\Sigma$.\\[4pt]
Độc lập trong chuẩn chung & $\mathbf X_1\perp\mathbf X_2\iff\Sigma_{12}=0$;
$\mathbf a'\mathbf X\perp\mathbf b'\mathbf X\iff\mathbf a'\Sigma\mathbf b=0$.\\[4pt]
Điều kiện, $\Sigma_{22}$ khả nghịch &
$\vmu_{1\mid2}=\vmu_1+\Sigma_{12}\Sigma_{22}^{-1}(\mathbf x_2-\vmu_2)$;
$\Sigma_{11\mid2}=\Sigma_{11}-\Sigma_{12}\Sigma_{22}^{-1}\Sigma_{21}$.\\[4pt]
Khoảng cách, tham số biết & $(\mathbf X-\vmu)'\Sigma^{-1}(\mathbf X-\vmu)\sim\chi^2_p$
 khi $\mathbf X$ chuẩn và $\Sigma$ xác định dương.\\[4pt]
MLE mẫu chuẩn, $A\succ0$ & $\widehat{\vmu}=\xbar$;
$\widehat\Sigma=A/n=(n-1)S/n$, với $A$ là tổng tích ngoài quanh $\xbar$.
Khác với $\Cov(\Xbar)=\Sigma/n$.\\[4pt]
Mẫu iid chuẩn & $\Xbar\sim\N_p(\vmu,\Sigma/n)$;
$(n-1)S\sim W_{p,n-1}(\Sigma)$; $\Xbar\perp S$.\\[4pt]
Mẫu iid, moment hữu hạn & $\sqrt n(\Xbar-\vmu)\xrightarrow{d}\N_p(0,\Sigma)$;
$S\xrightarrow{P}\Sigma$, với $p$ cố định.\\
\bottomrule
\end{longtable}

\subsection{Công thức suy diễn cho trung bình}
Đặt $\mathbf d=\xbar-\vmu_0$, $T^2=n\mathbf d'S^{-1}\mathbf d$.
\begin{longtable}{L{.28\linewidth}L{.63\linewidth}}
\toprule
Điều kiện / câu hỏi & Công thức\\
\midrule
\endhead
Mẫu iid chuẩn, $n>p$, $\Sigma\succ0$ &
$\dfrac{n-p}{p(n-1)}T^2\sim F_{p,n-p}$ dưới $H_0$;
$c_T^2=\dfrac{p(n-1)}{n-p}F_{p,n-p;1-\alpha}$.\\[5pt]
Miền tin cậy chính xác & $n(\xbar-\vmu)'S^{-1}(\xbar-\vmu)\le c_T^2$.\\[4pt]
Bonferroni thành phần & $\mathcal R_B=\prod_{i=1}^p
[\bar x_i-h_i;\bar x_i+h_i]$, $h_i=t_{n-1;1-\alpha/(2m)}\sqrt{s_{ii}/n}$.
Với danh sách chỉ các thành phần thì $m=p$;
miền không dùng hiệp phương sai chéo.\\[4pt]
Mọi tổ hợp, mẫu chuẩn & $\mathbf a'\xbar\pm\sqrt{c_T^2}\sqrt{\mathbf a'S\mathbf a/n}$.\\[4pt]
$m$ tổ hợp định trước, mẫu chuẩn & $\mathbf a_j'\xbar\pm
t_{n-1;1-\alpha/(2m)}\sqrt{\mathbf a_j'S\mathbf a_j/n}$;
phủ đồng thời ít nhất $1-\alpha$.\\[4pt]
Hiệu hai thành phần cùng đối tượng & Ước lượng $\bar x_i-\bar x_k$;
sai số chuẩn $\sqrt{(s_{ii}+s_{kk}-2s_{ik})/n}$.\\[4pt]
Mẫu lớn, iid, moment hữu hạn, $p$ cố định, $\Sigma\succ0$ & Thay $c_T^2$ bằng
$\chi^2_{p;1-\alpha}$; Bonferroni dùng $z_{1-\alpha/(2m)}$.
Độ tin cậy và mức kiểm định là xấp xỉ.\\[4pt]
Liên hệ với hợp lý, mẫu chuẩn & $\lambda_W=(1+T^2/(n-1))^{-1}$;
$\Lambda_{\mathrm{LR}}=\lambda_W^{n/2}$;
$T^2=(n-1)(|\widehat\Sigma_0|/|\widehat\Sigma|-1)$.\\
\bottomrule
\end{longtable}

\subsection{Từ điển ký hiệu ngắn}
\begin{center}
\begin{tabular}{lp{.70\linewidth}}
\toprule
Ký hiệu & Cách đọc\\
\midrule
$n,p$ & Số đối tượng và số biến của mỗi đối tượng.\\
$\vmu,\Sigma$ & Trung bình và hiệp phương sai của tổng thể.\\
$\Xbar,\xbar$ & Trung bình mẫu ngẫu nhiên và giá trị đã quan sát.\\
$S,\widehat\Sigma$ & Hiệp phương sai mẫu không chệch và MLE hiệp phương sai.\\
$S/n,\Theta$ & Ước lượng hiệp phương sai trung bình mẫu; $\Theta=\Sigma^{-1}$ là ma trận độ chính xác.\\
$\tr(B)$ & Tổng các phần tử đường chéo của ma trận vuông $B$.\\
$\Sigma_{12}$ & Khối hiệp phương sai chéo giữa hai nhóm biến.\\
$\perp$ & Độc lập ngẫu nhiên.\\
$\xrightarrow{P},\xrightarrow{d}$ & Hội tụ theo xác suất, hội tụ về phân phối.\\
$F_{r,s;\gamma},t_{\nu;\gamma}$ & Phân vị dưới mức $\gamma$; không phải đuôi trên $\gamma$.\\
$m$ & Số phát biểu trong danh sách Bonferroni định trước.\\
\bottomrule
\end{tabular}
\end{center}

\clearpage
\section{Bản đồ đối chiếu với giáo trình}
\label{sec:references}
Các vị trí dưới đây được đối chiếu với bản PDF địa phương của ấn bản thứ
sáu, Pearson New International Edition, bản quyền 2014.
Tên file năm 2013 không được dùng thay cho năm in ở trang bản quyền.
Trong Chương 4--5, muốn đến trang in 150, mở trang PDF 155; tương tự trang
in 235 là trang PDF 240.

\footnotesize
\begin{longtable}{L{.35\linewidth}L{.43\linewidth}r}
\toprule
Chủ đề & Vị trí trong sách & Trang in\\
\midrule
\endhead
Mật độ chuẩn nhiều chiều & Ch.~4, §4.2, (4-4) & 150\\
Hình học đường đồng mức & Ch.~4, §4.2, (4-7), Result 4.1 & 153\\
Tổ hợp tuyến tính một biến & Ch.~4, §4.2, Result 4.2 & 156\\
Biến đổi nhiều tổ hợp & Ch.~4, §4.2, Result 4.3 & 157\\
Phân phối biên & Ch.~4, §4.2, Result 4.4 & 158--159\\
Độc lập và hiệp phương sai 0 & Ch.~4, §4.2, Result 4.5 & 159--160\\
Phân phối có điều kiện & Ch.~4, §4.2, Result 4.6 & 160--161\\
Dạng hồi quy của trung bình điều kiện & Ch.~4, §4.2, (4-9) & 163\\
Mahalanobis và chi-bình phương & Ch.~4, §4.2, Result 4.7 & 163--164\\
Tách tổng tích ngoài & Ch.~4, §4.3, (4-14) & 169\\
Hàm hợp lý & Ch.~4, §4.3, (4-16), (4-17) & 170\\
Kiểm tra cực đại theo hiệp phương sai & Ch.~4, §4.3, Result 4.10 & 170--171\\
MLE trung bình, hiệp phương sai & Ch.~4, §4.3, Result 4.11 & 171--172\\
Wishart và phân phối mẫu & Ch.~4, §4.4, (4-22), (4-23) & 174\\
Luật số lớn & Ch.~4, §4.5, Result 4.12, (4-26), (4-27) & 175\\
CLT và dạng toàn phương giới hạn & Ch.~4, §4.5, Result 4.13, (4-28) & 176\\
Ôn kiểm định một chiều & Ch.~5, §5.2, (5-1)--(5-3) & 210--211\\
Thống kê Hotelling & Ch.~5, §5.2, (5-4) & 211\\
Phân phối $F$ và ngưỡng & Ch.~5, §5.2, (5-5)--(5-7) & 212\\
Bất biến & Ch.~5, §5.2, (5-9) và diễn giải & 215--216\\
Hợp lý cực đại dưới hai mô hình & Ch.~5, §5.3, (5-10)--(5-12) & 216--217\\
Wilks' lambda và $T^2$ & Ch.~5, §5.3, Result 5.1, (5-14) & 218\\
$T^2$ từ hai định thức & Ch.~5, §5.3, (5-15) & 219\\
Miền tin cậy elipsoid & Ch.~5, §5.4, (5-18), (5-19) & 221\\
Khoảng đồng thời mọi hướng & Ch.~5, §5.4, Result 5.3 & 225\\
Khoảng cho thành phần và hiệu & Ch.~5, §5.4, (5-24), (5-25) & 225--226\\
Bonferroni và khoảng thành phần & Ch.~5, §5.4, (5-28), (5-29) & 232\\
So sánh độ dài khoảng & Ch.~5, §5.4, (5-30) & 234\\
Suy diễn mẫu lớn & Ch.~5, §5.5, Results 5.4, 5.5, (5-31) & 235--236\\
\bottomrule
\end{longtable}
\normalsize

Một số vị trí trong bản tổng hợp ban đầu cần cập nhật khi dùng bản sách
này: phân phối có điều kiện vẫn thuộc §4.2; MLE hiệp phương sai thuộc
Result 4.11, không phải (4-14); công thức (5-15) ở trang 219; Bonferroni
(5-28), (5-29) ở trang 232. Nội dung giảng và các bảng ở đây dùng các
vị trí đã đối chiếu, đồng thời phân biệt rõ kết quả toán học với quy tắc
thực hành về kích thước mẫu.
